\exercisegroup

In all the Exercises for § 4, \(\mathfrak{C}\) denotes a quantified theory.

\begin{enumerate}
\def\labelenumi{\arabic{enumi}.}
\item
  Let A and B be relations in G, and let x be a letter which does not appear in A. Then \((\forall x)(A \Longrightarrow B) \Longleftrightarrow (A \Longrightarrow (\forall x)B)\) is a theorem in G.
\item
  Let \(A\) and \(B\) be relations in \(\mathfrak{C}\) , and let \(x\) be a letter, distinct from the constants of \(\mathfrak{C}\) , which does not appear in \(A\) . If \(B \Longrightarrow A\) is a theorem in \(\mathfrak{C}\) , then \(((\exists x)B) \Longrightarrow A\) is a theorem in \(\mathfrak{C}\) .
\item
  Let \(A\) be a relation in \(\mathfrak{C}\) and let \(x\) and \(y\) be letters. The relations
\end{enumerate}

\[
(\forall x) (\forall y) A \Longrightarrow (\forall x) ((x | y) A), \quad (\exists x) ((x | y) A) \Longrightarrow (\exists x) (\exists y) A
\]

are then theorems in \(\mathfrak{C}\) .

\begin{enumerate}
\def\labelenumi{\arabic{enumi}.}
\setcounter{enumi}{3}
\tightlist
\item
  Let A and B be relations in G and let x be a letter. Show that the relations
\end{enumerate}

\[
\begin{array}{l} (\forall x) (A \text { or } B) \Longrightarrow ((\forall x) A \text { or } (\exists x) B), \\ (\exists x) A \text { and } (\forall x) B \Longrightarrow (\exists x) (A \text { and } B) \end{array}
\]

are theorems in \(\mathfrak{C}\) .

\begin{enumerate}
\def\labelenumi{\arabic{enumi}.}
\setcounter{enumi}{4}
\tightlist
\item
  Let A and B be relations in G, and let x and y be letters. If x does not appear in B, nor y in A, then
\end{enumerate}

\[
(\forall x) (\forall y) (A \text {   and   } B) \iff ((\forall x) A \text {   and   } (\forall y) B)
\]

is a theorem in \(\mathfrak{C}\) .

\begin{enumerate}
\def\labelenumi{\arabic{enumi}.}
\setcounter{enumi}{5}
\item
  Let A and R be relations in G, and let x be a letter. Then the relations \((\exists_{A}x)R \Longrightarrow (\exists x)R\) , \((\forall x)R \Longrightarrow (\forall_{A}x)R\) are theorems in G.
\item
  Let \(A\) and \(R\) be relations in \(\mathfrak{C}\) , and let \(x\) be a letter distinct from the constants of \(\mathfrak{C}\) . If \(R \Longrightarrow A\) is a theorem in \(\mathfrak{C}\) , then
\end{enumerate}

\[
(\exists x) R \Longleftrightarrow (\exists_ {A} x) R
\]

is a theorem in \(\mathfrak{C}\) . If (not \(R) \Longrightarrow A\) is a theorem in \(\mathfrak{C}\) , then

\[
(\forall x) R \Longleftrightarrow (\forall_ {A} x) R
\]

is a theorem in \(\mathfrak{C}\) . In particular, if \(A\) is a theorem in \(\mathfrak{C}\) ,

\[
(\exists x) R \Longleftrightarrow (\exists_ {A} x) R \quad \text { and } \quad (\forall x) R \Longleftrightarrow (\forall_ {A} x) R
\]

are theorems in \(\mathfrak{C}\) .

\begin{enumerate}
\def\labelenumi{\arabic{enumi}.}
\setcounter{enumi}{7}
\tightlist
\item
  Let A and R be relations in C, let T be a term in C and let x be a letter. If \((T|x)A\) is a theorem in C, then \((T|x)R \Longrightarrow (\exists_{A}x)R\) and \((\forall_{A}x)R \Longrightarrow (T|x)R\) are theorems in C.
\end{enumerate}

